\documentclass{article}
\usepackage[utf8]{inputenc}
\usepackage{hyperref}
\usepackage{url}
\title{Yamazaki?? ??????????????? Tang????? ?????????? ?????? ?????-?????? ????? ????????????, ???????????? ???????? ????}
\author{sci-adk (deterministic render)}
\date{\today}

\begin{document}
\maketitle

\noindent\textit{Draft compiled by sci-adk from Spec \texttt{pfas-rice-oos-lipid} (v1). Belief state is revisable as Evidence accrues.}

\section{Goal}
Yamazaki?? ??????????????? Tang????? ?????????? ?????? ?????-?????? ????? ????????????, ???????????? ???????? ????
?????? ????????? Tang 2026 ???????? TF?? ??????? ???????? ?????????? (??????????? ????????????? ??????? RMSE??
in-sample ???????? ?????(\textasciitilde{}0.52)?? ??????????, ????? PFOS ??????????? ??????????).

\section{Background}
??????? ??????(`runs/pfas-rice-oos-tang`)??? **???????????? ?????**(????? monotone f\_xy, Tang ????????)??
?????? ????????? Tang 2026?? ???????? ???????????? **????????? ??????**??? ???????(??????? log10 RMSE 1.232
vs in-sample ???????? 0.519). ??????? ??????? **PFOS**(??-K\_PL ?????????)??, ???????????? ???????
\textasciitilde{}40--200$\times$ ??????????????. ?????? ????? sub-investigation(`runs/pfas-rice-longchain`, LC1/LC2 SUPPORTED)???
**B-??????? ?????-?????? ?????? ????? ???**(g\_xy$\cdot$C, g\_ph$\cdot$C, K\_PL-gated)?? ?????????????? in-sample ?????
(Yamazaki C10--C12) ???????? ????????? ????????, ????? Chen 2025(???--??? K\_MW?? ???????????? ?????? --- ????? ?????
???????$\cdot$???????????????? ??????)?? ?????? ?????? corroborate?????. **????????????? `LIPID\_LOADING` ????????
Yamazaki(PFDoDA ????)?? ????????? ??????? Tang?? ????????? ????? ????????**, ?? ???????????? Tang?? ???????????
????? ??????? ???????? ??????? ?????.

\section{Method}
`model\_api.tang\_tf\_validation`(ORYZA2000 biomass, dry-weight TF, f\_xy\_source="recommended"=?????
monotone$\cdot$Tang ????????)???? PFOA$\cdot$PFOS$\cdot$GenX 3????? ???? **lipid\_loading=False(???????????? ????????)** ??
**lipid\_loading=True(Yamazaki-?????? g\_xy/g\_ph)** ?? ??????? ????? TF??? Tang ?????? TF(dw, 0.1 $\mu$g/g)??
?????????(`validation/oos\_tang\_lipid.py`). ?? ??????? ??????? log10 RMSE??? in-sample ???????? ???????
?????????. lipid ???????? Yamazaki ??????????????? Tang?? ???? ????? ?????????? ????????? ????????. ??????
(Tang 2026)?? ????? ???????.

Planned approaches:
\begin{itemize}
  \item `model\_api.tang\_tf\_validation`(ORYZA2000 biomass, dry-weight TF, f\_xy\_source="recommended"=?????
monotone$\cdot$Tang ????????)???? PFOA$\cdot$PFOS$\cdot$GenX 3????? ???? **lipid\_loading=False(???????????? ????????)** ??
**lipid\_loading=True(Yamazaki-?????? g\_xy/g\_ph)** ?? ??????? ????? TF??? Tang ?????? TF(dw, 0.1 $\mu$g/g)??
?????????(`validation/oos\_tang\_lipid.py`)
  \item ?? ??????? ??????? log10 RMSE??? in-sample ???????? ???????
?????????
  \item lipid ???????? Yamazaki ??????????????? Tang?? ???? ????? ?????????? ????????? ????????
  \item ??????
(Tang 2026)?? ????? ???????.
\end{itemize}

\section{Hypotheses and findings}

\subsection{Yamazaki?? ??????????????? Tang????? ?????????? ?????? ?????-?????? ????? ????????????, ???????????? ???????? ????
?????? ????????? Tang 2026 ???????? TF?? ??????? ???????? ?????????? (??????????? ????????????? ??????? RMSE??
in-sample ???????? ?????(\textasciitilde{}0.52)?? ??????????, ????? PFOS ??????????? ??????????).}
\begin{itemize}
  \item Hypothesis id: \texttt{hyp-001} (exploratory)
  \item Decision rule (qualitative): Expert judgment based on evidence
  \item \textbf{Status: supported} --- confidence strong (graded)
  \item Basis: qualitative judge: Out-of-sample generalization SUCCEEDS: the lipid-facilitated loading mechanism (K\_PL-gated; constants fit on Yamazaki excl. PFDoDA, NOT on Tang) drops the independent-dataset Tang OOS log10 RMSE from 1.232 to 0.516 -- matching the in-sample Tang-refit (0.519) without fitting anything to Tang. It fixes the dominant free-anion failure (PFOS \textasciitilde{}40-200x under -> spot-on stalk), consistent with the membrane-partition mechanism Chen 2025 independently confirms. This is the project's first strong cross-dataset out-of-sample predictive success and resolves the hyp-oos-tang REFUTED baseline at the mechanism level. Residual GenX ether over-prediction (provisional offset, separate from lipid loading) and a PFOS-endosperm under-prediction remain, so 'supports' (mechanism generalizes), not a perfect fit. Hypothesis supported.
  \item \textit{Evidence validity:} referent=empirical; data\_source(s)=measured
\end{itemize}

\section{Evidence}
\begin{itemize}
  \item \texttt{evi-oos-lipid} (experiment\_run): point=0.516, finding=Out-of-sample GENERALIZATION on the INDEPENDENT Tang 2026 dataset. The K\_PL-gated lipid-facilitated loading mechanism (L
\end{itemize}

\section{Pending agent judgments}
These hypotheses use a non-numeric rule and await an in-session agent verdict (no autonomous LLM call):
\begin{itemize}
  \item \texttt{hyp-001} (qualitative): Expert judgment based on evidence
  \item finding: Out-of-sample GENERALIZATION on the INDEPENDENT Tang 2026 dataset. The K\_PL-gated lipid-facilitated loading mechanism (LIPID\_LOADING constants fit on YAMAZAKI excl. PFDoDA, NOT on Tang -- see docs/fxy\_longchain\_lipid\_exploration.md) drops the OOS log10 RMSE from 1.232 (free-anion baseline) to 0.516 -- matching the in-sample Tang-refit level (0.519) WITHOUT touching any parameter for Tang. The dominant baseline error (PFOS, the high-K\_PL sulfonate) is largely fixed: stalk 0.013->0.620 vs Tang 0.571, endosperm 0.004->0.194 vs 0.944, leaf 0.039->1.536 vs 0.684 (the \textasciitilde{}40-200x under collapses). PFOA also improves (stalk 0.485->1.378 vs 2.222; endosperm 0.136->0.444 vs 1.371). Residual (honestly recorded): GenX (ether) stays OVER-predicted (stalk 8.3->5.6, leaf 17.7->11.6 vs \textasciitilde{}1) -- the ether f\_xy offset is provisional/condition-dependent, NOT a lipid-loading failure; and PFOS endosperm is still \textasciitilde{}5x under. The mechanism is corroborated by Chen 2025 (membrane K\_MW monotone with chain length; the lipid pool carries the most-sorptive species). This is the project's first strong cross-dataset out-of-sample predictive success: a mechanism fit on one dataset predicts an independent dataset's per-organ pattern.
\end{itemize}

\section{References}
No literature cited.

\end{document}